%!TEX program = pdflatex
% ============================================================
%  TEMPLATE CV « Informatique x Recherche » — une page
%  Alexis MIALON
%  Compilation : pdflatex cv_alexis_mialon.tex  (Overleaf OK)
%  Aucune police ni paquet exotique.
%
%  Structure : Profil > Expériences (industrie + recherche)
%              > Open source & projets > Formation > Compétences
%  Pour ajouter une ligne : copier un bloc \entry{...}{...}{...}
%  suivi d'un environnement itemize.
% ============================================================
\documentclass[10pt,a4paper]{article}

\usepackage[utf8]{inputenc}
\usepackage[T1]{fontenc}
\usepackage[scaled=0.92]{helvet}
\renewcommand{\familydefault}{\sfdefault}
\usepackage{microtype}
\usepackage[top=0.9cm,bottom=0.8cm,left=1.4cm,right=1.4cm]{geometry}
\usepackage{xcolor}
\usepackage{titlesec}
\usepackage[shortlabels]{enumitem}
\usepackage[hidelinks,pdftitle={CV - Alexis Mialon},pdfauthor={Alexis Mialon}]{hyperref}

% ---------- Couleurs (modifier ici pour changer le thème) ----------
\definecolor{accent}{HTML}{1F6FB2}
\definecolor{muted}{HTML}{5A5A5A}

% ---------- Mise en page ----------
\pagestyle{empty}
\setlength{\parindent}{0pt}
\setlength{\parskip}{0pt}
\setlist[itemize]{leftmargin=1.0em,itemsep=0pt,topsep=1pt,parsep=0pt,label=\textcolor{accent}{\textbullet}}

\titleformat{\section}{\normalsize\bfseries\color{accent}}{}{0pt}{\MakeUppercase}[\vspace{-4.5pt}\textcolor{accent}{\rule{\linewidth}{0.5pt}}]
\titlespacing*{\section}{0pt}{5.5pt}{2.5pt}

% ---------- Commandes ----------
% \entry{Titre}{Organisation / contexte}{Date}
\newcommand{\entry}[3]{%
  \vspace{2.5pt}%
  \textbf{#1}\ \textcolor{muted}{\small|}\ \textcolor{accent}{\small #2}\hfill\textcolor{muted}{\small #3}\par}
\newcommand{\tag}[1]{\textcolor{muted}{\small #1}}

\begin{document}
\fontsize{9.6}{11.6}\selectfont

% ---------- En-tête ----------
\begin{center}
  {\fontsize{21}{24}\selectfont\bfseries Alexis MIALON}\\[2pt]
  {\normalsize\color{accent}\bfseries Élève-ingénieur EPITA \,·\, IA agentique \,·\, C++ \,·\, Cryptanalyse}\\[3pt]
  {\footnotesize\color{muted}
    Lyon \,·\, \href{mailto:alexis.mialon69@gmail.com}{alexis.mialon69@gmail.com} \,·\, 06\,76\,15\,68\,88 \,·\,
    \href{https://www.linkedin.com/in/alexis-mialon-79117b329/}{linkedin.com/in/alexis-mialon-79117b329}\\[1pt]
    \href{https://github.com/Vorpalin}{github.com/Vorpalin} \,·\,
    \href{https://vorpalin.github.io/}{vorpalin.github.io}}
\end{center}

\vspace{-5pt}
% ---------- Profil ----------
\section{Profil}
Élève-ingénieur à l'EPITA (diplôme prévu en 2028), stagiaire IA chez Opensee sur une solution d'IA agentique
pour l'analyse de données financières et étudiant-chercheur en cryptanalyse différentielle (C++ parallèle).
Contributeur open source (deux PR fusionnées dans \texttt{langchain-aws}). Je cherche à travailler là où
l'ingénierie logicielle, l'IA et la recherche se rencontrent.

% ---------- Expériences ----------
\section{Expériences}

\entry{Stagiaire IA}{Opensee, Puteaux}{Sept. 2026 -- présent}
\begin{itemize}
  \item Contribution au développement d'\textbf{Agensee}, solution d'IA agentique dédiée à l'analyse de données financières (équipe IA).
\end{itemize}

\entry{Étudiant-chercheur, équipe Sécurité \& Systèmes}{Laboratoire de Recherche de l'EPITA, Lyon}{Depuis fév. 2026}
\begin{itemize}
  \item Optimisation de la minimisation des \textbf{S-boxes actives} dans la recherche de caractéristiques différentielles de l'outil de cryptanalyse \textbf{TAGADA} : algorithmes de \textbf{programmation dynamique multithreadés en C++}, adaptés aux contraintes de la bibliothèque.
  \item Automatisation de la récupération des tours de chiffrement (parsing JSON), benchmarks conteneurisés (Docker), veille scientifique et synthèse des avancées à l'équipe.
\end{itemize}

\entry{Assistant pédagogique en programmation}{EPITA (cycle préparatoire) et SupBiotech, Lyon}{2025 -- 2026}
\begin{itemize}
  \item Accompagnement quotidien d'étudiants de première année en \textbf{C} et \textbf{Python}, Git et Linux (NixOS) ; relecture des sujets de TP et des tests, correction (Criterion, GitLab).
\end{itemize}

% ---------- Open source & projets ----------
\section{Open source \& projets}

\entry{Contributions à \texttt{langchain-aws} (LangChain)}{2 PR fusionnées : \href{https://github.com/langchain-ai/langchain-aws/pull/1389}{\#1389}, \href{https://github.com/langchain-ai/langchain-aws/pull/1415}{\#1415}}{Oct. 2026}
\begin{itemize}
  \item Ajout des niveaux d'effort de raisonnement (\texttt{reasoning\_effort\_levels}) aux profils de modèles Amazon Bedrock ; valeurs vérifiées à la source, lint et tests validés.
\end{itemize}

\entry{Emulators Collection}{C++20, WebAssembly, React \;·\; \href{https://github.com/Vorpalin/emulators-collection}{GitHub}}{Projet personnel}
\begin{itemize}
  \item Émulateurs \textbf{CHIP-8, Super-CHIP, Atari 2600 et Game Boy} : cœurs C++20 indépendants du navigateur, compilés en \textbf{WebAssembly} (Emscripten), frontend React/TypeScript/Vite ; CI GitHub Actions, analyse de sécurité, releases automatiques, Docker.
\end{itemize}

\entry{PING, IDE en ligne pour débutants}{Java, React ; équipe agile de 5, EPITA}{Depuis mars 2026}
\begin{itemize}
  \item Conception produit (état de l'art, interviews, maquettes), interface React, tests backend automatisés par pipeline CI/CD GitLab, documentation Javadoc.
\end{itemize}

\entry{42sh, shell POSIX en C}{équipe agile de 4, EPITA}{Janv. 2026}
\begin{itemize}
  \item Lexer et parser, architecture fondée sur les patrons \textbf{visiteur} et \textbf{composite}, revue des merge requests GitLab. \tag{Autre projet : Yakamon, API REST Java en couches avec JUnit (1 semaine).}
\end{itemize}

% ---------- Formation ----------
\section{Formation \& distinctions}

\entry{Diplôme d'ingénieur en informatique}{EPITA, Lyon}{Sept. 2023 -- 2028}
\entry{Semestre à l'étranger, Computer Science}{Centria University of Applied Sciences, Finlande}{Jan. -- mai 2025}

\vspace{2pt}
\textbf{Top Scorer}, Challenge du Meilleur étudiant en informatique de France 2026 \,·\,
\textbf{TOEIC} 960/990 \,·\, Baccalauréat général (2023)

% ---------- Compétences ----------
\section{Compétences}
\begin{tabular}{@{}p{2.6cm}p{14.4cm}@{}}
  \textbf{Langages}        & C, C++ (C++20), Python, Java, C\#, JavaScript / TypeScript, SQL (PostgreSQL, MySQL) \\[1pt]
  \textbf{IA \& recherche} & IA agentique, cryptanalyse différentielle, programmation dynamique, calcul parallèle \\[1pt]
  \textbf{Web \& WASM}     & React, Vite, WebAssembly (Emscripten), HTML/CSS, API REST \\[1pt]
  \textbf{Outils}         & Git (GitHub, GitLab), CI/CD, Docker, CMake, Linux (NixOS), JUnit, Criterion \\[1pt]
  \textbf{Langues}         & Français (natif), Anglais (C1, TOEIC 960/990), Italien (notions)
\end{tabular}

% ---------- Divers ----------
\section{Divers}
\textbf{Tennis de table} : Départemental~1 par équipes, pré-régional en individuel \,·\, Cinéma \,·\,
Bénévole (2016--2023) : animation d'activités, référent de stand (équipe d'une vingtaine de personnes).

\end{document}
